\documentclass{article}
\usepackage[utf8]{inputenc}
\usepackage[brazil]{babel}

\title{Título do Artigo}
\author{Autor}
\date{\today}

\begin{document}
Nesse trabalho nosso grupo analisou dados relacionados à saúde mental utilizando técnicas de Machine Learning para identificar padrões associados à ansiedade. O foco principal do estudo foi entender a relação entre a qualidade e a duração do sono e os níveis de ansiedade, verificando como essas informações podem contribuir para a criação de modelos preditivos.

Para o desenvolvimento do projeto, inicialmente foram importadas as bibliotecas necessárias. Utilizou-se o Pandas para leitura e manipulação do conjunto de dados, o Matplotlib para a construção de gráficos e o Scikit-learn para o treinamento e a avaliação do modelo de Machine Learning.

Em seguida, o arquivo CSV foi carregado e analisado utilizando as funções head(), info() e describe(), assim permitindo visualizar as primeiras linhas do dataset, verificou os tipos de dados presentes e obter um resumo estatístico das variáveis numéricas. Também foi realizada uma análise estatística mais detalhada, considerando média, mediana, desvio padrão, valores mínimos e máximos, o que possibilitou uma melhor compreensão da distribuição dos dados.


Os resultados foram obtidos por meio da avaliação do desempenho, que  foi feita pelas métricas de acurácia, precisão, recall e F1-Score, além da análise da matriz de confusão e do relatório de classificação. Também foi analisada a importância das variáveis para identificar quais características tiveram maior influência nas decisões do modelo, destacando especialmente os relacionados ao sono.

Por fim, foram realizados três experimentos adicionais, alterando a quantidade de árvores da Random Forest, a profundidade máxima das árvores e a proporção entre os dados de treinamento e teste. Os resultados obtidos foram satisfatórios em todas as configurações avaliadas, demonstrando que o modelo apresentou bom desempenho e estabilidade.

\maketitle

\end{document}
